\section{Conclusion}

This paper has rewritten freeform Fourier-operator metasurface design as a response-centered inverse problem defined over wavelength, angle, and polarization. On that basis, it develops a coherent physical loop comprising fabrication-aware structure generation, simulation-based labeling, score-aware data preparation, forward surrogate training, conditional diffusion sampling, and simulation-based reranking. The key scientific point is that the optical target should be specified in the same response space in which the final device is verified, placing the present study in closer contact with objective-driven photonic inverse design and response-faithful electromagnetic modeling \cite{an2019deep,ma2019probabilistic,phillips2021ensemble}.

Within this formulation, a second-order differentiator with large NA, high transmission, and broadband operation reaches a leading level of verified performance across multiple wavelengths, while polarization-independent, polarization-multiplexed, fourth-order, low-pass, and spatiotemporal targets can likewise be handled at a competitive level within the same unified inference framework. The broader contribution is the demonstration that multiple operator families can be described, generated, and judged within one physically consistent design language.

The broader implication is methodological. Freeform metasurface inverse design is better treated as target-conditioned feasible-set discovery under electromagnetic verification than as single-geometry regression or task-label prediction detached from the transfer function. The present study provides one concrete realization of that principle, and it points naturally toward future extensions involving richer polarization control, full complex-response targets, broader physical design spaces, tighter fabrication models, and larger-scale physics-guided datasets.
